% -*- coding: utf-8 -*-
% !TEX program = xelatex
\documentclass[plain,noanswer,medmath]{../../template/jnuexam}
\usepackage{../../template/kysx}

\begin{document}


\examtitle{1989}{1}


\makepart{填空题}{本题满分15分，每小题3分}


\begin{problem}
已知$f'(3)=2$，则 $\lim_{h\to 0}\,\frac{f(3-h)-f(3)}{2h}=$\fillin{}.
\end{problem}


\begin{problem}
设$f(x)$是连续函数，且$f(x)=x+2\int_0^1f(t) \dt$，则$f(x)=$\fillin{}.
\end{problem}


\begin{problem}
设平面曲线$L$为下半圆周$y=-\sqrt{1-x^2}$，则曲线积分$\int_L (x^2+y^2)\ds=$\fillin{}.
\end{problem}


\begin{problem}
向量场$\vec{u}(x,y,z)=xy^2\vi+y\e^z\vj+x\ln(1+z^2)\vk$
在点$P(1,1,0)$处的散度$\div\vec{u}=$\fillin{}.
\end{problem}


\begin{problem}
设矩阵$A=\begin{pmatrix}
  3 & 0 & 0 \\
  1 & 4 & 0 \\
  0 & 0 & 3
\end{pmatrix}$, $E=\begin{pmatrix}
  1 & 0 & 0 \\
  0 & 1 & 0 \\
  0 & 0 & 1
\end{pmatrix}$，则逆矩阵$(A-2E)^{-1}=$\fillin{}.
\end{problem}


\makepart{选择题}{本题满分15分，每小题3分}


\begin{problem}
当$x>0$时，曲线$y=x\sin\frac{1}{x}$\pickout{}
\begin{abcd}
\item 有且仅有水平渐近线．
\item 有且仅有铅直渐近线．
\item 既有水平渐近线，也有铅直渐近线．
\item 既无水平渐近线，也无铅直渐近线．
\end{abcd}
\end{problem}


\begin{problem}
已知曲面$z=4-x^2-y^2$上点$P$处的切平面平行于平面$2x+2y+z-1=0$，则点$P$的坐标是\pickout{}
\begin{abcd}
\item $(1,-1,2)$．
\item $(-1,1,2)$．
\item $(1,1,2)$．
\item $(-1,-1,2)$．
\end{abcd}
\end{problem}


\begin{problem}
设线性无关的函数$y_1$, $y_2$, $y_3$都是二阶非齐次线性方程
$y''+p(x) y'+q(x)y=f(x)$的解，$C_1$, $C_2$是任意常数，则该非齐次方程的通解是\pickout{}
\begin{abcd}
\item $C_1y_1+C_2y_2+y_3$．
\item $C_1y_1+C_2y_2-(C_1+C_2)y_3$．
\item $C_1y_1+C_2y_2-(1-C_1-C_2)y_3$．
\item $C_1y_1+C_2y_2+(1-C_1-C_2)y_3$．
\end{abcd}
\end{problem}


\begin{problem}
设函数$f(x)=x^2$, $0\le x<1$，
而$\smash[b]{S(x)=\sum_{n=1}^{\infty}b_n\sin n\pi x}$, $-\infty <x<+\infty$，
其中$b_n=2\int_0^1f(x)\sin n\pi x\dx$, $n=1,2,3,\cdots$，
则$S\left(-\frac{1}{2}\right)$等于\pickout{}
\begin{abcd}
\item $-\frac{1}{2}$．
\item $-\frac{1}{4}$．
\item $\frac{1}{4}$．
\item $\frac{1}{2}$．
\end{abcd}
\end{problem}


\begin{problem}
设$A$是$n$阶矩阵，且$A$的行列式$|A|=0$，则$A$中\pickout{}
\begin{abcd}
\item 必有一列元素全为$0$．
\item 必有两列元素对应成比例．
\item 必有一列向量是其余列向量的线性组合．
\item 任一列向量是其余列向量的线性组合．
\end{abcd}
\end{problem}


\makepart{计算题}{本题满分15分，每小题5分}


\begin{problem}
设$z=f(2x-y)+g(x,xy)$，其中函数$f(t)$二阶可导，$g(u,v)$具有连续的二阶偏导数，
求$\frac{\pd^2 z}{\pdx\pd y}$．
\end{problem}


\begin{problem}
设曲线积分$\int_C xy^2 \dx+y\varphi(x)\dy$与路径无关，
其中$\varphi (x)$具有连续的导数，且$\varphi (0)=0$,
计算$\int_{(0,0)}^{(1,1)}xy^2 \dx+y\varphi(x)\dy$的值．
\end{problem}


\begin{problem}
计算三重积分$\iiint_{\Omega}(x+z)\d V$，其中$\Omega$
是由曲面$z=\sqrt{x^2+y^2}$与$z=\sqrt{1-x^2-y^2}$所围成的区域．
\end{problem}


\makepart{}{本题满分6分}

\begin{problem}
将函数$f(x)=\arctan\frac{1+x}{1-x}$展为$x$的幂级数．
\end{problem}


\makepart{}{本题满分7分}

\begin{problem}
设$f(x)=\sin x-\int_0^x (x-t)f(t)\dt$，其中$f$为连续函数，求$f(x)$．
\end{problem}


\makepart{}{本题满分7分}

\begin{problem}
证明方程$\ln x=\frac{x}{\e}-\int_0^{\pi}\!\!\!\sqrt{1-\cos 2x}\dx$
在区间$(0,+\infty)$内有且仅有两个不同实根．
\end{problem}


\makepart{}{本题满分6分}

\begin{problem}
问$\lambda$为何值时，线性方程组
$$\begin{cases}
  x_1+ x_3=\lambda \\
  4x_1+x_2+2x_3=\lambda +2 \\
  6x_1+x_2+4x_3=2\lambda +3
\end{cases}$$
有解，并求出解的一般形式．
\end{problem}


\makepart{}{本题满分8分}

\begin{problem}
假设$\lambda$为$n$阶可逆矩阵$A$的一个特征值，证明：\par
\begin{enumerate*}
\item $\frac{1}{\lambda}$为$A^{-1}$的特征值；
\item $\frac{|A|}{\lambda}$为$A$的伴随矩阵$A^*$的特征值．
\end{enumerate*}
\end{problem}


\makepart{}{本题满分9分}

\begin{problem}
设半径为$R$的球面$\Sigma$的球心在定球面$x^2+y^2+z^2=a^2(a>0)$上，
问当$R$为何值时，球面$\Sigma$在定球面内部的那部分的面积最大？
\end{problem}


\makepart{填空题}{本题满分6分，每小题2分}


\begin{problem}
已知随机事件$A$的概率$P(A)=0.5$，随机事件$B$的概率$P(B)=0.6$，及条件概率
$P(B|A)=0.8$，则和事件$A\cup B$的概率$P(A\cup B)=$\fillin{}.
\end{problem}


\begin{problem}
甲、乙两人独立地对同一目标射击一次，其命中率分别为$0.6$和$0.5$．
现已知目标被命中，则它是甲射中的概率为\fillin{}.
\end{problem}


\begin{problem}
若随机变量$\xi$在$(1,6)$上服从均匀分布，则方程$x^2+\xi x+1=0$有实根的概率是\fillin{}.
\end{problem}


\makepart{}{本题满分6分}

\begin{problem}
设随机变量$X$与$Y$独立，且$X$服从均值为$1$、标准差(均方差)为$\sqrt{2}$的正态分布，
而$Y$服从标准正态分布．试求随机变量$Z=2X-Y+3$的概率密度函数．
\end{problem}


\end{document}
